\scanpage{151}
la composition connue des classes de correspondance.
\inserted{On a catégorie additive (objet nul est $\varnothing$, somme est $\amalg$).}

Si on admet que l'équiv. hom. $=$ équiv. numérique, alors on déduit (par th. de finitude de la cohomologie) que les $\operatorname{Fl}(X,Y)$ sont des \unsure{esp.} vectoriels de dim finie sur $\mathbf Q$ ; si on admet l'inégalité de Hodge, \unsure{un|les} $\operatorname{Fl}(X,X)$ devient une algèbre involutive du type de \unsure{Castelnuovo}-Weil.... donc est semi-simple. D'où une catégorie \deleted{de} $\mathcal M$, mais on \unsure{va} considérer celle associée à $\mathcal V^\circ$, de sorte que on a un foncteur \uline{contravariant} additif universel
\diagram{\begin{tikzcd}
\mathcal V^\circ \arrow[r,dash,"\uline h"] & \mathcal M
\end{tikzcd}}

D'ailleurs, on a un foncteur évident
\[
H^{*}(,\mathbf Q_\ell):\mathcal V^\circ\longrightarrow\operatorname{Mod}_{\mathrm{t.f.}}(\mathbf Q_\ell)
\]
en associant à tout $X\in\operatorname{Ob}\mathcal V$ le module $H^{*}(X,\mathbf Q_\ell)$. Ce foncteur est \uline{fidèle}, par définition \unsure{même} des $\operatorname{Hom}$ dans $\mathcal V$ ($+$ dualité ...).

On trouve un foncteur exact fidèle \unsure{can.}
\[
T_\ell:\mathcal M\unsure{\mathbin{-}|\longrightarrow}\operatorname{Mod}_{\mathrm{t.f.}}(\mathbf Q_\ell)
\]
appelé \uline{foncteur de Tate}.
\mspage{150}
